\section{Normal forms, witnesses and refutations}
\label{sec:normal}

This section connects three kinds of evidence. A finite truth table can supply witnesses for a Skolem expansion. A ground truth table can supply a resolution refutation. A cut-free derivation can supply the very ground instances to which the second construction applies. All transformations below take numerical codes as input. Statements involving evaluation in a structure use a coded nonempty finite domain $[d]=\{0,\ldots,d-1\}$ unless expressly stated otherwise.

\subsection{Prenex and clausal normal forms}

Use $\neg,\wedge,\vee,\to,\forall,\exists$, together with atomic equality, relation atoms and truth constants. Treat $A\leftrightarrow B$ as an abbreviation for $(A\to B)\wedge(B\to A)$. Formula sizes below count the expanded trees, including term nodes. In a traversal, give each binder a distinct temporary name greater than every free identifier. For de Bruijn input, the name of a bound occurrence is found by looking up its binder in the traversal stack. A reverse traversal with a scope stack returns de Bruijn syntax. Thus a prefix can be moved without ever changing which binder owns an occurrence.

\begin{theorem}[Polarity and prenex conversion]
There are primitive-recursive transformations $N^+,N^-,P$. The first two yield formulas in negation normal form and evaluate respectively to the original bit and its complement. $P$ yields a string of quantifiers followed by a quantifier-free matrix with the same free variables and the same evaluation as the input.
\end{theorem}
\begin{proof}
A positive atom stays an atom and a negative atom becomes its negation. A negation exchanges the two polarities. On a conjunction the positive branch joins the positive children by $\wedge$, and the negative branch joins the negative children by $\vee$. On a disjunction interchange $\wedge$ and $\vee$ in those instructions. For implication the positive output is $N^-(A)\vee N^+(B)$ and the negative output is $N^+(A)\wedge N^-(B)$. The positive quantifier branch keeps the quantifier, and the negative branch interchanges $\forall,\exists$ and takes the negative body. Complement truth constants at negative polarity. The two outputs have the stated bits at an atom. The four possibilities for a pair of child bits prove each binary step. Complementing a finite conjunction of bits gives the disjunction of their complements: the empty folds have bits $1,0$, and appending one bit applies the binary De Morgan law. The dual identity is identical with the two bits exchanged. These identities prove the quantifier steps.

Run a second traversal on the positive output. A literal returns the empty prefix and itself. A quantifier prepends its binder to the child's prefix. Given prefix--matrix pairs $(U,B),(V,C)$ for two children joined by $\star\in\{\wedge,\vee\}$, return $(UV,B\star C)$. For the proof, move the binders of $U$ across the right child and then those of $V$ across the left matrix, from outside to inside. At each move one operand has a fixed bit $b$ independent of the variable moved. Both finite folds commute with $z\mapsto z\star b$: if this map is the identity this is immediate, and if it is constant, a fold of $d\ge1$ copies of that constant is the same constant. Distinct fresh names give exactly the independence required. This also explains the domain hypothesis: on the empty domain $\top\vee\exists x\,\bot$ and $\exists x(\top\vee\bot)$ have different bits.

Postorder processing constructs at most two polarity outputs per node. A single polarity output visits each input occurrence once and adds at most one negation per atomic occurrence. Prenex conversion moves binder records and retains the matrix nodes. The temporary-name range has length the number of binders. For the induction verifying conversion, fix the input and structure, and conjoin the two desired bit identities over the processed subtrees and all assignments to names appearing there. If $v$ names appear, there are at most $d^v$ assignments. Each induction assertion is one primitive-recursive bit test.
\end{proof}

The same fold calculation gives the article's identities $\neg\forall x A\equiv\exists x\neg A$ and $\neg\exists x A\equiv\forall x\neg A$ in these models. Two like quantifiers commute: a double conjunction is $1$ exactly when every ordered-pair entry is $1$. For two disjunctions, scan for the first $1$ and transpose its pair of indices. Universal quantification distributes over conjunction and existential quantification over disjunction by distributing the corresponding finite folds entry by entry. Freshness is essential in the rule for pulling a quantifier across an operand. The preceding proof records the point where it is used.

\begin{theorem}[Skolem tables with an explicit choice function]
Let $F=Q_1x_1\cdots Q_qx_q B$ be a closed prenex sentence with quantifier-free $B$. For every existential position $i$ add a fresh function $s_i$ of arity the number $u_i$ of preceding universal variables, in their prefix order. Substitute $s_i(\bar x_{<i}^{\forall})$ for $x_i$ in $B$, keep the universal prefix, and call the result $S$. A primitive-recursive function expands every finite model of $F$ to a model of $S$ on the same domain. Every finite model of $S$ reduces to a model of $F$.
\end{theorem}
\begin{proof}
For each prefix history $h$ of length $i$, compute the bit $T_i(h)$ of the remaining suffix. Fill the last row with the matrix bits, then compute the preceding rows by finite AND or OR according to the next quantifier. Now fill the function tables in prefix order. At an existential position and a fixed tuple of preceding universal moves, reconstruct the intervening existential moves using tables already filled. This gives one history $h$. Record the least $a<d$ for which the next row reads $T_i(h,a)=1$ (here $i$ is the one-based position just played). Use $0$ when there is no such entry. This fills the entire table even when the root bit is zero. If the root bit is one, a universal move preserves a one for each choice and an existential table selects a one whenever its parent is one. Induction on the play length therefore makes the final matrix one for every universal tuple. Evaluation after substitution gives exactly this matrix value.

Conversely, suppose all universal tuples give one under some supplied Skolem tables. Every history consistent with the tables has all of its complete continuations true when later existential moves also follow the tables. Work backward from length $q$. At a universal position apply the preceding assertion to every child. At an existential position the table entry gives one child to which it applies. The suffix bit is one in both cases, hence the root bit is one. Each row assertion ranges over at most $d^q$ histories. The suffix array has $\sum_{j=0}^q d^j$ entries, and the new function tables use $\sum_{i:Q_i=\exists}d^{u_i}$ entries. All bounds include the constant-function case $d^0=1$. Each replacement of a variable occurrence adds at most $q$ argument nodes. These counts give primitive recursions on prefix positions and bounded scans for every operation.
\end{proof}

\begin{proposition}[Two clause conversions]
A quantifier-free matrix in negation normal form has an equivalent conjunction of clauses. If it has $m$ nodes, one may use at most $2^m$ clauses of width at most $m$. Alternatively, after adding fresh relation symbols, at most $3m+1$ clauses suffice and every satisfying interpretation of the matrix at all universal tuples has a unique expansion satisfying all these clauses.
\end{proposition}
\begin{proof}
Use one singleton clause for a literal, no clauses for $\top$, and the single empty clause for $\bot$. Concatenate clause lists at a conjunction. At a disjunction form every pairwise union of one clause from each child list. If both children have a false clause, their union is false. A bounded scan finds such clauses whenever both lists have Boolean value zero. If either list is true, every union is true. This proves the disjunction rule including empty lists. Conjunction adds list lengths, disjunction multiplies them, and a union's width is bounded by the sum of widths. Induction gives the two bounds.

For the second construction let $\bar u$ be all the universal variables. Use the literal itself as its output wire. For each binary node introduce a distinct predicate $p(\bar u)$. If its children's wires are $a,b$, impose the following clauses, where a bar complements a literal:
\[
\begin{array}{c|lll}
p\leftrightarrow(a\wedge b)&\bar p\vee a&\bar p\vee b&p\vee\bar a\vee\bar b\\
p\leftrightarrow(a\vee b)&p\vee\bar a&p\vee\bar b&\bar p\vee a\vee b .
\end{array}
\]
A constant node uses a fresh wire fixed by one unit clause, and a final unit asserts the root wire. For each tuple the three clauses force exactly the displayed Boolean gate value, as checking the two possible output bits shows. Consequently postorder evaluation supplies the only possible new predicate entries. The root clause holds for all tuples exactly when the matrix does. Arity zero covers a sentence with no universal variables. The bound counts clause occurrences, not the total length of their argument lists.
\end{proof}

\subsection{Ground resolution and analytic Herbrand witnesses}

A clause is a finite set of signed atoms. Repetition is immaterial. A list of clauses is true when every clause contains a true literal. Its empty clause is always false. A ground atom is an atom whose terms have no variables. In propositional tests distinct ground atoms are treated as independent bits until explicit equality or congruence conditions are added.

\begin{theorem}[Bounded resolution completeness and soundness]
For any finite family $G$ of ground clauses on $k$ distinct atoms, primitive-recursive functions return either a satisfying assignment or a resolution refutation with at most $2^{k+1}-1$ nodes. Each leaf names an input clause. Each internal node names a pivot and the union of its parents after deleting their designated opposite pivot literals.
\end{theorem}
\begin{proof}
Enumerate the $2^k$ assignments and return the least satisfying one if it exists. Otherwise at every full assignment choose the first input clause it falsifies. For a partial assignment $s$ to the first $j$ atoms let $B_s$ contain the $j$ signed literals falsified by those choices. At a leaf the selected clause is a subset of $B_s$. Work up the complete binary tree, retaining at each node a derived clause $D\subseteq B_s$. If $p$ is its next atom, the children have clauses $D_0\subseteq B_s\cup\{p\}$ and $D_1\subseteq B_s\cup\{\neg p\}$. If either lacks its displayed pivot, retain that child alone. Otherwise resolve the children on $p$. The union of the two remainders is a subset of $B_s$. At the root this forces the empty clause. Pruning children only decreases the given node bound. The induction assertion is the conjunction of the subset tests in one tree level.

For soundness fix a valuation satisfying both parents of a resolution step. If the pivot bit is zero, the parent containing its positive occurrence must have a true remainder. If the bit is one, the parent containing its negative occurrence must have a true remainder. In either case the union is true. Induction over a proof therefore precludes an empty conclusion under any satisfying valuation. Both the checker and the witness construction are bounded traversals.
\end{proof}

For universally closed clauses, an additional sound rule instantiates a clause by any displayed well-scoped substitution. Its truth follows by interpreting the substituted terms, then using the assignment so obtained in the original clause. A first-order resolution certificate may supply substitutions on two separately renamed clauses and a check that the resulting pivot atoms coincide. It may also factor two literals which that substitution makes equal. Literal-set union implements the latter deletion. These operations require only syntactic comparison. Soundness follows from the preceding rule and the ground pivot argument applied at each assignment. In particular a finite, propositionally inconsistent family of ground instances refutes its universal parent clauses. For example, $R(a)$, $\neg R(x)\vee R(f(x))$ and $\neg R(f(f(a)))$ have a four-leaf certificate. Use instances $x=a,f(a)$ of the middle clause and resolve successively to obtain $R(f(a))$, $R(f(f(a)))$, and the empty clause.

\begin{theorem}[Two-way cut-free Herbrand extraction]
\label{thm:herbrand}
In equality-free classical one-sided sequent calculus, take a sequent of closed formulas $E_i=\exists\bar x_i\,A_i$, with quantifier-free matrices in negation normal form. From an explicit cut-free tree proof one can extract a finite list of ground instances of the $A_i$ whose disjunction is a propositional tautology. Conversely, such a list and a successful tautology check give a cut-free proof of the original sequent. Both maps are primitive recursive. If the language has no closed term, use a fresh parameter $c$ in the ground language.
\end{theorem}
\begin{proof}
Use multisets (or lists with exchange), weakening, contraction, the one-sided rules for conjunction, disjunction, and existential instantiation, and initial sequents containing complementary literals or $\top$. The conjunction rule has two premises with the same context. The disjunction rule has one premise containing both children. The existential rule takes $\Gamma,A[t/x]$ to $\Gamma,\exists x A$. These are precisely the logical rules needed on paths read backward from this endsequent: weakening removes a formula, contraction duplicates one, and the logical rules peel its prefix or its matrix. Equivalently, restrict an analytic calculus in negation normal form to this endsequent. Atomic initial sequents and explicit structural rules fix the proof format.

Replace all free proof parameters by a fixed closed term (adjoin $c$ if necessary), after avoiding capture. Trace each root occurrence backward with a tag recording its root index and witnesses already supplied in prefix order. Copy tags at contraction and at branching. Whenever a prefix is exhausted, record its ground matrix and tuple, including occurrences whose prefix was empty at the root. Further decomposition of its matrix requires no new record. There are at most the number of root occurrences plus the number of existential rule nodes many records.

Fix any bit assignment on the recorded atoms, extended by zero on any other proof atoms. Suppose every recorded matrix is false. Give every still-prefixed occurrence the temporary value false. Give all quantifier-free occurrences their usual propositional value. Starting from the root, whose occurrences are all false, follow a path with every occurrence false. At an existential step the exposed occurrence still has a nonempty prefix or is a just-recorded false matrix. At a disjunction both children are false. At a conjunction choose a false child and follow its premise. Backward weakening, contraction and exchange preserve the assertion. No initial sequent can have all occurrences false. This contradicts the finite path ending at a leaf. Since the purported failure was a bit assignment, a loop over all assignments proves the tautology test succeeds. This argument follows only bounded paths. It also works when the recorded list is empty, in which case no proof could have supplied it.

Conversely start with the quantifier-free sequent consisting of the listed instances. Decompose a disjunction into its two disjuncts in one premise, or a conjunction into its two alternative premises. Propositional validity of the conclusion implies validity of each resulting premise by the two gate truth tables. Continue until $\top$ closes a branch or only literals and $\bot$ remain. A valid leaf without $\top$ must contain opposite literals, since otherwise assigning each occurring positive atom zero and each occurring negative atom one falsifies the whole leaf. This constructs a cut-free derivation using at most $2^{L+1}-1$ tree positions for $L$ connective occurrences on the initial list: along a branch the count strictly decreases. Now restore the recorded existential binders around each matrix in reverse prefix order using the recorded terms. Contract repeated copies of each $E_i$ and weaken in any missing $E_i$. If $c$ was adjoined and was absent from the intended signature, replace it throughout the proof by a single fresh free proof variable. No free variable enters the endsequent.

For numerical verification, keep tags and instance lists on a bounded traversal of the supplied tree. The invariant for extraction is indexed by proof nodes and valuations of the finite set of proof atoms. Grounding the free parameters and instantiating displayed terms are syntax traversals. The reverse construction recurses on the connective count with the displayed exponential bound. This gives quantifier-free induction for the advertised maps.
\end{proof}

Taking $A_i$ to be the conjunction of complements of a universal clause gives a direct corollary. Cut-free refutability of a finite equality-free universal clause family is equivalent to the existence of an explicit finite list of ground instances having a resolution refutation. The forward map uses Theorem~\ref{thm:herbrand} and bounded resolution. The reverse map uses resolution soundness, finite truth-table duality and the reverse construction in Theorem~\ref{thm:herbrand}. A witness list supplies its own term-depth bound. No bound on search depth as a function of the clause family is invoked.

\subsection{A decision with two verifiable outputs}

\begin{theorem}[Bernays--Sch\"onfinkel--Ramsey with equality]
\label{thm:bsr}
Let $F=\exists x_1\cdots x_e\,\forall y_1\cdots y_u\,B$ be closed. The finite signature has relation symbols, equality, $b$ constants, and no other functions. There is a primitive-recursive decision algorithm which returns either a model of $F$ of size at most $h=\max(1,b+e)$, or a checkable classical first-order derivation of $F\to\bot$. In particular the same bit decides whether $F$ has \emph{any} coded finite model and detects infeasibility even when equality can identify different existential witnesses or constants.
\end{theorem}
\begin{proof}
Use a name for each original constant and each existential witness. Use one dummy name only if there are no names. There are $h$ names in total. Let $G$ be the propositional CNF obtained by the following instructions.
\begin{enumerate}
\item Give the CNF conversion of $B$ the witness names for $\bar x$ and instantiate $\bar y$ in all $h^u$ ways by names. Include every resulting clause.
\item Represent each equality of names $i,j$ by an atom $E_{ij}$, and each relation on name tuples by an atom $R_{\bar i}$. Include the units $E_{ii}$, all clauses expressing $E_{ij}\to E_{ji}$, and all clauses expressing $E_{ij}\wedge E_{jk}\to E_{ik}$.
\item For each relation $R$ of arity $r$ and each pair of name tuples include the clause for
\[
  \left(\bigwedge_{t=1}^r E_{i_t j_t}\right)\wedge R_{\bar i}
  \ \longrightarrow\ R_{\bar j}.
\]
\end{enumerate}
There are at most $k=h^2+\sum_R h^{\operatorname{arity}(R)}$ atoms. In addition to the ground matrix clauses, the construction uses at most $h+h^2+h^3+\sum_R h^{2\operatorname{arity}(R)}$ equality and relation clauses. Apply bounded resolution.

From a satisfying assignment, form the equivalence classes of $E$ using as representative the least equivalent name. Reflexivity guarantees a successful bounded search. Symmetry and transitivity imply two names have the same representative exactly when $E$ relates them. Interpret names by their representatives, and assign $R$ on representative tuples using the assigned bits. The congruence clauses in both tuple orders imply that replacing every coordinate by its representative preserves its bit. Thus for every tuple of name representatives the matrix evaluates to one. Every representative is itself a name, so all universal assignments are covered. The designated witness names supply the original existential moves. Ordering the representatives and replacing each by its rank gives a genuine model on an initial segment of size between one and $h$.

For the opposite branch, a resolution leaf has one of two justifications. For an equality or congruence leaf use logical equality: repeatedly replace one tuple coordinate using its equality, thereby obtaining the displayed congruence implication. Reflexivity is an axiom, symmetry follows by replacing in $i=i$, and transitivity by replacing in $i=j$. For a matrix leaf, assume $F$, open its existential witnesses using fresh parameters, and apply universal elimination at the indicated tuple of names. The inferred ground matrix entails each of its CNF clauses. To obtain classical proof codes uniformly for this last step, split on each of the finitely many atoms by excluded middle. Under a complete assignment, formulas with value one have proofs, and formulas with value zero have proofs of their negation, from the signed atomic assumptions: recurse over the formula, using the connective rules and the known signs of its children. In the true-matrix case the indicated clause has value one. In the false-matrix case the assumed matrix gives a contradiction. Disjunction elimination on the successive excluded middles discharges the signed assumptions. The same method gives the small propositional conversions from the equality implications to clauses. A resolution step is derivable by case analysis on both parent disjunctions: if either supplies a remainder literal introduce it into the conclusion. The opposite pivot pair gives a contradiction and hence the conclusion. Repetition and order in a clause can be changed by the same disjunction rules. The empty-clause root therefore gives $\bot$. Discharge the fresh witness parameters by existential elimination in reverse order, then discharge $F$. The dummy name may be translated to a fresh free variable, as in Theorem~\ref{thm:herbrand}. Classical rules used here occur in the \emph{output proof}. Its generation and checking use bounded operations in iPRA.

Finally, any given finite model of $F$ supplies satisfying bits for $G$. Read its constants, scan for an existential tuple making the universal suffix true, and interpret the dummy name by element zero when needed. Record the equalities and relations of these elements. Every matrix instance is true, and actual equality supplies all congruence clauses. Hence the resolution branch is impossible for this model. The three constructions show exhaustiveness and exclusivity constructively. In ordinary external semantics the same finite restriction applies to any inhabited model with actual equality by taking the interpretations of these finitely many names. The internal result already decides between a finite model and a formal refutation.
\end{proof}

\paragraph{An equality test that changes the answer.}
For $F=\exists x\exists y\,(x=y\wedge P(x)\wedge\neg P(y))$ the ground matrix forces $E_{12},P_1,\neg P_2$. The congruence clause $\neg E_{12}\vee\neg P_1\vee P_2$ resolves against these three units to give the empty clause. For $\exists x\exists y\,(x=y\wedge P(x))$ both witnesses may name the same element and the quotient returns a one-element model. Assigning distinct domain elements to fresh constants by default would lose that possibility.

\paragraph{Lean implementation order.}
Use \texttt{Nat} identifiers and carry the prefix as a list of pairs (quantifier tag, fresh variable identifier). Define \texttt{polar : Bool $\to$ Formula $\to$ Formula}, \texttt{splitPrefix} returning the pair (prefix, matrix), and \texttt{suffixRows} by primitive recursion on the reversed prefix. Prove the bit equations for finite folds before proving correctness of \texttt{polar} and \texttt{splitPrefix}. A Skolem-table row takes a preceding universal tuple to the least accepted child of the reconstructed history. Its forward invariant is truth at every history generated from that tuple. Its reverse invariant is truth of the original suffix whenever all consistent completions satisfy the matrix. For resolution use signed integers or pairs (atom index, sign), a sorted deduplicated list for a clause, and an inductive certificate with constructors for indexed input leaves and binary pivots. The induction lemma returns both an accepted certificate and a clause contained in the literals falsified by the current prefix. For Herbrand extraction store a root index and a list of witnesses on each sequent occurrence, copying that record at contraction. Formalize the false-premise selector as a function of a valued sequent, its last rule and the valuation. Its iteration is bounded by tree height. For the fragment decision, build the lists of ground relation and equality atoms by tuple enumeration and prove representative equality in both directions before the structural induction on the ground matrix. Each correctness lemma is transferred to code arithmetic by collecting checks over the finite stage, assignments, and tuples specified in its proof.
